\documentclass{article}
\begin{document}
\title{A dense point-cloud flag can survive clipping and defeat the finite-point filter}
\maketitle

\paragraph{Compared interface.}
The pinned RGB-D point path sets \texttt{is\_dense=true} before
applying the native mask. It can replace all three coordinates of
a point by infinity; it does not revise that flag. The selected
CoSTAR launch forwards camera points through a bridge, and the
pinned ROS bridge binary-copy conversion preserves the flag and
the serialized coordinates. The reference downstream chain here
uses pinned melodic-era \texttt{pcl\_conversions} and PCL 1.8.1.
These comparator versions and a consumer using them are explicit
fixture choices, not a finding that the deployed robot ran this
filter or these exact package builds.

\paragraph{Native exercise.}
Inject three raw points in each selected frame. Move one target
among all three positions while its raw value is finite near,
finite beyond the software depth clip, positive or negative
infinity, or NaN. Include both depth-subscriber modes. Use selected
native sensor control assignment, callback copy, clip/mask and
packed XYZ, native ROS conversion of data/metadata, PCL conversion
and typed cloud unpacking, and PCL's selected finite filter and
extrema routines. Stand-ins supply frame buffers, message types,
known standard XYZ field mapping and scalar vector operations;
no wire serialization timing or rendered target is exercised.
The PCL point-cloud type documents a dense flag as a promise that
all point data are finite; the fixture focuses on XYZ.

The stock setter marks every case dense. In all 21 cases with an
invalid XYZ point, the flag remains true at the PCL input.
The native filter bypasses the finite predicates when this
flag is true and returns every point and index. In the 15
cases involving infinities, the selected extent calculation
also returns a nonfinite X extrema value in the fixture. No
NaN-specific extent claim is made about another library build.

\paragraph{What the repair does and does not supply.}
Demoting the unchecked flag before invoking the original PCL
filter makes the existing finite check run. It produces a
finite output cloud from the known valid points in all cases
and returns their original point indices. In particular a
compacted position is not the same index as the original ray
once a previous point was dropped. A direct finite mask and
index map agree exactly. This repairs the downstream
finiteness contract, not visibility or acquisition. It cannot
reconstruct the discarded far observation, distinguish clipping
from missing rendering from the final message alone, or certify
free space where no valid observation remains. The producer-side
scope in the parent study is still required for its far cue.
No independent benchmark, field measurement, physical safety or
architecture acceptance is claimed.
\end{document}
